\chapter{Eigenvalues, singular values, and principal components}\label{ch:pca}
\begin{readergoal}Understand the spectrum of a matrix, its geometry, and the exact optimization solved by PCA.\end{readergoal}
\section{Eigenvectors as preserved directions}
For a square matrix $A$, an eigenvector $v\ne0$ satisfies $Av=\lambda v$. The scalar $\lambda$ describes the action on that direction. A real symmetric matrix admits an orthonormal basis of real eigenvectors and a factorization $A=Q\Lambda Q^\top$. A covariance matrix is symmetric positive semidefinite, so its eigenvalues are nonnegative.

\section{Singular value decomposition}
Every real matrix $X\in\mathbb R^{T\times N}$ admits an SVD $X=U\Sigma V^\top$ with nonnegative singular values in $\Sigma$ and orthonormal singular vectors. This construction exists even for rectangular or rank-deficient matrices. The right singular vectors describe directions in variable space; the left singular vectors describe directions in observation space. Truncating to the first $k$ components gives the best rank-$k$ approximation in Frobenius norm (Eckart--Young theorem), provided singular values are ordered descending.

\section{Deriving PCA}
Let $X_c$ denote the column-centered observations and $S=X_c^\top X_c/(T-1)$. A projected observation onto a unit vector $v$ has sample variance $v^\top Sv$. The maximum-variance direction solves
\[v_1=\arg\max_{\|v\|=1}v^\top Sv.\]
Using the Lagrangian $\mathcal L(v,\lambda)=v^\top Sv-\lambda(v^\top v-1)$ gives $Sv=\lambda v$. The maximizing direction is an eigenvector of the largest eigenvalue. Subsequent components maximize variance subject to orthogonality. If $X_c=U\Sigma V^\top$, then $\lambda_j=\sigma_j^2/(T-1)$.

\section{What PCA does not promise}
PCA maximizes sample variance, not class separation, behavioral relevance, temporal predictability, or a causal mechanism. Large variance can come from artifacts, gain fluctuations, or unmodeled experimental factors. Cross-validation, null models, and careful centering are essential. As \cite{CunninghamYu2014,Churchland2012} discuss, structured population dynamics may require objectives different from variance maximization.

\begin{workedexample}[An effective dimension]
If $\lambda_1,\ldots,\lambda_N$ are nonnegative variances, the participation ratio $d_{\rm PR}=(\sum_i\lambda_i)^2/\sum_i\lambda_i^2$ equals $N$ when all eigenvalues are equal and equals one when only one is positive. It is a continuous summary of spectral concentration, not an integer count of ``real'' sources.\end{workedexample}
\begin{exercises}
\item Derive $\lambda_j=\sigma_j^2/(T-1)$ for centered observations.
\item Compare participation ratios for $(10,0,0)$ and $(4,3,3)$.
\item Explain why replacing trials by a trial average can remove single-trial fluctuations.
\item Give a case where the leading principal component tracks recording artifact rather than neural activity.
\end{exercises}